\begin{frame}
\titlepage
\note{Welcome the class and name the question for today: when does a tidy tree help, and when does it hide a relationship? Tell students that every example uses synthetic data only.}
\end{frame}

\begin{frame}{Week 12 continuity: the model comes first}
\princ{MODEL VS PROCEDURE}
\begin{block}{Carry this sentence forward}
\large A correct procedure on the wrong model is still wrong.
\end{block}
\begin{itemize}
\item Last week, a one-way street meant the two directions had different edges.
\item This week, a tree will make a tidy picture of relationships.
\item Before tracing the picture, ask what one edge means.
\end{itemize}
\predict{If the arrows mean dependency but the tree edges mean containment, can the same picture answer both questions?}
\note{Connect to Week 12 in one minute. A procedure can be perfectly executed and still answer the wrong question if its model is wrong. Today the trap is a change in edge meaning: dependency is a work-order relationship; containment is a displayed-under relationship. Ask for a prediction before showing the graph.}
\end{frame}

\begin{frame}{Today: model, trace, audit}
\princ{CHECK THE WORK}
\begin{enumerate}
\item Start with a concrete seven-item project graph.
\item Turn it into a directory-style tree, and name the edge meaning.
\item Trace the tree three ways: preorder, postorder, and breadth-first.
\item Audit what the tree dropped; then test a decision tree's limit.
\end{enumerate}
\begin{columns}[T]
\begin{column}{0.48\textwidth}
\predict{Which trace will visit both top-level branches before their children?}
\end{column}
\begin{column}{0.48\textwidth}
\begin{alertblock}{Slide legend}
\textcolor{blue!60!black}{PREDICT} = guess first;\\
\textcolor{green!40!black}{What you should see} = visible result;\\
\textcolor{red!60!black}{Common mistake} = failure mode.
\end{alertblock}
\end{column}
\end{columns}
\note{Preview the sequence and the loop. Students should predict before code runs, then compare the visible result, then explain the gap. Do not start with abstract definitions; the project picture comes next. The answer to the prediction is breadth-first, but do not give it away yet.}
\end{frame}

\begin{frame}{Words we will use}
\princ{MODEL VS PROCEDURE}
\small
\begin{tabular}{@{}p{3.0cm}p{9.9cm}@{}}
\textbf{graph} & A set of items and relationships between them.\\[2pt]
\textbf{vertex} & One item in a graph; the dot in a picture.\\[2pt]
\textbf{directed edge} & An arrow from one vertex to another; direction matters.\\[2pt]
\textbf{dependency} & A relationship saying one item must be completed before another.\\[2pt]
\textbf{tree} & A rooted, connected structure with one parent for each non-root item.\\[2pt]
\textbf{containment} & A display relationship: one item is shown under another.\\[2pt]
\textbf{traversal} & A rule for the order in which we visit vertices.\\[2pt]
\textbf{assert} & A test statement that fails when a claim is false.\\[2pt]
\textbf{decision tree} & A tree whose branches represent choices and outcomes.\\
\end{tabular}
\note{Do not read every row aloud. Point out that graph, directed edge, tree, and traversal are the words that will recur. Emphasize the two meanings: dependency tells us about completion; containment tells us about display. If a term is new later, return to this slide.}
\end{frame}
